\section{Introduction}
A finite causal experiment has two distinct memory problems. A simulator
must retain enough of the past to produce the required future law. An
auditor must retain enough of its observations to select and execute a
statistically useful continuation. A restriction at the end of training
does not, by itself, restrict the computation that produced that state.
The question addressed here is how a decision continuation bound extends
to a machine whose register is charged throughout acquisition,
randomization and validation.

We work with declared causal interfaces and actual marked probability
laws. Candidate behavior images need not be convex. Preparation count,
retained labels, transition precision and the schedule clock are distinct
resources. An external clock chooses a predetermined transition; it does
not store sample history. An autonomous realization must retain that
clock as part of its state. This distinction is essential even for a
Bernoulli experiment.

For a compact candidate image $B\subset\Delta(\Omega)$ and target $Q$,
a reset auditor observes independent trials from one fixed $b\in B$,
then selects an event before one fresh candidate and one fresh target
validation. Its signed score is the target indicator minus the candidate
indicator. A randomized event has a response $h\in[0,1]^\Omega$. Define
\[
 \mathfrak V_K(B,Q)=\max_{h_1,\ldots,h_K}
                    \min_{b\in B}\max_{i\le K}(Q-b)h_i,
 \qquad
 \kappa_c(B,Q)=\min\{K:\mathfrak V_K(B,Q)>c\}.
\]
The decision-cut classification and its finite-sample approximation are
proved in Section~\ref{sec:v20-spectrum}. The new point is that a
response cover can be learned and implemented without an upstream
histogram of unbounded width.

\begin{theorem}[Multicut realization and certification]
\label{thm:v21-main}
The following statements hold under the interfaces and resource
conventions just described.

\smallskip\noindent
\textup{(i)} For $d=|\Omega|$, let $W_c$ be the minimum peak width over
all finite training schedules, including every random-bit and validation
cut, and let $W_c^{\rm rat}$ restrict to finite fair-bit implementations.
Then
\[
 \kappa_c(B,Q)\le W_c(B,Q)\le W_c^{\rm rat}(B,Q)
                                \le12d\,\kappa_c(B,Q).
\]
For a $K$-response cover, a finite-horizon tournament realizes this bound
with $2LR(K-1)$ training observations, where
$L=\lceil\log(2/\delta)/(4\theta)\rceil$ and
$R=\lceil2^L\log(2/\delta)\rceil$. Its uniform score loss is at most
$2^{1-r}+2\rho+4(K-1)\theta+2(K-1)\delta$, with dyadic precision $r$
and response-mean calibration error $\rho$.

\smallskip\noindent
\textup{(ii)} Approximate normalized continuation lists with compatible
one-step updates and local causal errors $\eta_t$ realize their entire
width profile with total error at most $\sum_t\eta_t$. Full-profile
risk bounds also charge the sum of the per-call calibration errors and
the row-rounding errors. Composition multiplies the simultaneous
simulator and auditor profiles. A shared word-capacity constraint turns
weighted residual obstructions at several cuts into a single Bayesian
regret lower bound. A family of point-source experiments has exact
decision spectra and matching linear peak-width growth; overlapping
independent retention tasks have simultaneous sharp exact profiles and
quantitative approximate lower bounds.

\smallskip\noindent
\textup{(iii)} In the prepared two-sphere collision family, for
$9/10\le d\le11/10$ and $0<\sigma\le1/24$, the row-indexed decision
complexity for score greater than $2/5$ is exactly three. A deterministic
physical audit has peak width at most twelve throughout report-bit
acquisition and validation, with an explicit finite preparation count.
The earlier $399$-training-trial audit is also retained, with peak width
$40602$ and its exact simultaneous residual profile. For the complete
physical task these give $3\le W_{2/5}^{\rm phys}\le12$; the upper
endpoint is not claimed to be an exact integer optimum.

\smallskip\noindent
\textup{(iv)} Independent repetition turns either positive expected-score
certificate into a test with both errors at most $\alpha$, charging the
score accumulator and every preparation. The short physical audit needs
$J\ge50\log(1/\alpha)$ repetitions. The physical certificate and the
bounded common-preparation Gaussian/entropic conditional-law consumer
are stable under the moving-collision perturbations specified below.
\end{theorem}
\begin{proof}
Part (i) is Theorem~\ref{thm:v21-endtoend}. Part (ii) follows from
Theorems~\ref{thm:v21-compatible}, \ref{thm:v21-transfer},
\ref{thm:v21-joint}, \ref{thm:v21-point}, \ref{thm:v21-intervals} and
Proposition~\ref{prop:v21-compose}. Part (iii) is
Theorem~\ref{thm:v21-physical} together with the preserved
Theorems~\ref{thm:v20-profile} and \ref{thm:v20-physical}. Part (iv)
follows from Theorems~\ref{thm:v21-confidence} and
\ref{thm:v21-fixedwitness}, Corollary~\ref{cor:v21-regimes}, and the
common-preparation transport results in
Theorems~\ref{thm:v18-microscopic} and \ref{thm:v19-global}.
\end{proof}

The central bound is about approximate statistical value over all finite
training lengths. It is not the exact residual profile of a prescribed
word function. Conversely, the exact profile of a particular selector
cannot be used as an all-algorithm approximate-memory lower bound. The
joint residual theorem explains when wordwise distinctions do carry
statistical weight, through the actual likelihoods and Bayes margins.
The growing point-source and retention families show that the
full-profile statements have nonconstant resource consequences.

The physical application makes the training tradeoff especially visible.
The short majority audit uses a moderate number of preparations and a
large counter quotient. The second audit uses a finite rare-block
comparison, a register of four bits, and a very long schedule. Neither
is presented as optimizing preparation count and peak width
simultaneously. Charging the clock changes the latter bound by the
explicit autonomous simulation formula. Likewise, perturbing a source
used in training generally introduces a training-count term; the
source-preserving moving-collision estimate is identified separately.

\subsection{Classical ingredients and the contribution}
Blackwell--Le Cam comparison and earlier filtered comparison provide the
statistical background~\cite{Blackwell1953,LeCam1964,Torgersen1991,
Norberg2002,Weisshaupt2002}. In Weisshaupt's report, Lemma 1 and
Theorem 5 on pp.~14--16 use compact convex classes of filtered
stochastic operators. Here the bounded register is a causal realization
constraint; its behavior image is not replaced by a convex class of
unpriced randomized machines. The full original Norberg article has not
been obtained for a theorem-by-theorem priority audit. No claim that its
theorems omit a particular feature is inferred from that limitation.

Positive-cone realization and residual minimization are classical
mechanisms~\cite{Fliess1975,Nerode1958}. Cover's Theorem 2, pp.~833--834,
already gives four-state clocked Bernoulli threshold learning, and his
Section 4 explicitly discusses the external clock~\cite{Cover1969}.
Our use of finite blocks is an implementation of that mechanism, not a
claim to have discovered finite-state learning. The work here is the
quantitative passage from a general nonconvex decision cover to a
complete rational training-and-validation machine, its all-cut resource
and perturbation accounting, and its integration with distribution-sensitive
lower bounds and a microscopic marked experiment.

\subsection{Organization and the broader program}
Sections~\ref{sec:v21-precision}--\ref{sec:v21-joint} develop the new
realization and obstruction results after the underlying comparison and
decision spectra. The finite-reset, streaming, nonlinear-testing,
physical-scattering and conditional-transport arguments are retained in
full. Sections~\ref{sec:v21-physical} and \ref{sec:v21-confidence}
complete the physical resource and statistical-confidence consequences.
Appendix~\ref{sec:v20-program} gives the actual proof dependencies for
all eleven historical components. The independent A2 geometric chain,
the nonlinear kinetic B4 targets and the broader C2 operator targets
retain their own proof obligations. This finite controlled resource
layer contributes to the program's general resource--precision--error
objective without replacing those model theorems by nominal arrows.
